%==============================================================================
% Chapter 3 — Classical ML & Model Evaluation
% Curriculum: Phase 2, Chapter 3 of docs/AI_ML_Master_Checklist.md
% Companion code: projects/03_classical_ml/
%
% Section skeleton only. Figures for this chapter go in theory/figures/.
%==============================================================================
\chapter{Data, Classical ML \& Model Evaluation}
\label{ch:classical-ml}

The first applied chapter, and the one that sets the pattern for the rest of
Phase~2: derive the model, then build it.\projectref{projects/03\_classical\_ml/} The models here are the ones that
still win on tabular data, and their mathematics is transparent enough to work
through by hand --- ordinary least squares has a closed form, and the gradient
boosting objective is a second-order Taylor expansion of exactly the kind seen
in \cref{eq:taylor}.

The chapter opens on data rather than on models, because that is where the time
actually goes and because a leak introduced during preprocessing invalidates
every number that follows. Everything here is \texttt{pandas} and
\texttt{polars}: the Compendium stays in one language end to end, so there is no
SQL in this chapter and no relational database in \cref{ch:capstone}.

It then pairs the models with the evaluation machinery they need. That pairing
is intentional: a model without a validation strategy is not a result, the
bias--variance decomposition explains \emph{why} regularisation works rather
than merely asserting that it does, and \cref{sec:statistical-inference} is what
turns "B scored higher" into a defensible claim. Everything in the last third of
the chapter applies unchanged to the deep models of
\cref{ch:deep-learning,ch:gen-ai}.

%==============================================================================
\section{Data Loading and Exploratory Analysis}
\label{sec:eda}

\subsection{Dataframes in \texttt{pandas} and \texttt{polars}}

\subsection{Vectorised Operations Instead of Loops}

\subsection{Joins, Merges and Grouped Aggregation}

\subsection{Profiling an Unfamiliar Dataset}

\subsection{Inspecting Distributions and Correlations}

%==============================================================================
\section{Data Cleaning}
\label{sec:data-cleaning}

\subsection{Missing-Data Mechanisms: MCAR, MAR and MNAR}

\subsection{Imputation Strategies}

\subsection{Outlier Detection}

\subsection{Duplicates, Dtype Coercion and Memory Footprint}

%==============================================================================
\section{Feature Engineering}
\label{sec:feature-engineering}

\subsection{Scaling and Standardisation}

\subsection{Categorical Encoding: One-Hot, Ordinal and Target}

\subsection{Interaction and Polynomial Features}

\subsection{Binning and Cyclical Encoding}

%==============================================================================
\section{Dimensionality Reduction}
\label{sec:dimensionality-reduction}

% PCA is derived from the covariance matrix in \cref{sec:unsupervised}; this
% section is about using it, and about the reduction methods that are only ever
% visualisation tools.

\subsection{PCA in Practice}

\subsection{t-SNE and UMAP for Visualisation}

\subsection{When Reduction Is the Wrong Choice}

%==============================================================================
\section{Linear Models}
\label{sec:linear-models}

\subsection{Ordinary Least Squares and the Normal Equations}

\subsection{Logistic Regression and the Sigmoid Map}

\subsection{Ridge and Lasso Regularisation}

%==============================================================================
\section{Tree-Based Models}
\label{sec:trees}

\subsection{Decision Trees: Information Gain and Gini Impurity}

\subsection{Random Forests and the Mathematics of Bagging}

\subsection{Gradient Boosting Machines}

\subsection{XGBoost and LightGBM: the Taylor Expansion of the Loss}

%==============================================================================
\section{Support Vector Machines}
\label{sec:svm}

\subsection{Max-Margin Classification}

\subsection{The Primal and Dual Problems}

\subsection{The Kernel Trick}

%==============================================================================
\section{Unsupervised Learning}
\label{sec:unsupervised}

\subsection{K-Means and Lloyd's Algorithm}

\subsection{Hierarchical Clustering}

\subsection{Principal Component Analysis}
\label{sec:pca}

% Derive PCA via the eigendecomposition of the covariance matrix, then connect
% it back to \cref{thm:svd}: the principal components are the right singular
% vectors of the centred data matrix, which is why PCA is usually computed by
% SVD rather than by forming the covariance matrix at all (see
% \cref{sec:numerical-stability}).

%==============================================================================
\section{Recommender Systems}
\label{sec:recommenders}

% The first direct application of \cref{eq:eckart-young} outside compression:
% a ratings matrix is mostly missing, and factorising it into user and item
% factors is a low-rank approximation of the observed entries.

\subsection{Collaborative versus Content-Based Filtering}

\subsection{Matrix Factorization}

\subsection{Implicit Feedback}

\subsection{The Cold-Start Problem}

\subsection{Ranking Metrics: MAP@k and NDCG}

%==============================================================================
\section{The Bias--Variance Tradeoff}
\label{sec:bias-variance}

\subsection{Decomposing the Expected Error}

\subsection{Reading Regularisation Through the Decomposition}

%==============================================================================
\section{Validation Strategies}
\label{sec:validation}

\subsection{K-Fold Cross-Validation}

\subsection{Stratified Sampling}

\subsection{Temporal Splits}

%==============================================================================
\section{Metrics}
\label{sec:metrics}

\subsection{The Confusion Matrix}

\subsection{Precision, Recall and the F1-Score}

\subsection{ROC and the Area Under the Curve}

\subsection{Regression Metrics: RMSE, MAE and \texorpdfstring{$R^2$}{R-squared}}

\subsection{Comparing Two Models Honestly}

% A gap in validation scores is an estimate with a standard error. Apply
% \cref{sec:statistical-inference}: confidence intervals on the difference,
% paired tests across folds, and correction when many models are compared.

%==============================================================================
\section{Hyperparameter Optimization}
\label{sec:hyperparameter-optimization}

\subsection{Grid and Random Search}

\subsection{Successive Halving and Early Stopping of Trials}

\subsection{Bayesian Optimization}

% The surrogate model is a Gaussian process, developed in
% \cref{sec:gaussian-processes}. State the acquisition-function idea here and
% defer the GP mathematics to \cref{ch:bayesian-ml}.

\subsection{\texttt{Optuna} in Practice}

\subsection{Nested Cross-Validation and Tuning Without Leaking}

%==============================================================================
\section{Model Interpretability}
\label{sec:interpretability}

% The expected accompaniment to any ensemble result in practice --- a gradient
% boosting model presented without SHAP values reads as unfinished work, so this
% section is core rather than optional.

\subsection{Why Built-In Tree Importances Mislead}

\subsection{Permutation Importance}

\subsection{Partial Dependence and ICE Plots}

\subsection{Shapley Values}

% From cooperative game theory: the unique attribution satisfying efficiency,
% symmetry, dummy and additivity. Derive the Shapley value as the average
% marginal contribution over all feature orderings, then show why the exact
% computation is exponential in the number of features.

\subsection{SHAP for Ensembles}

% TreeSHAP computes exact Shapley values for tree ensembles in polynomial time,
% which is why SHAP became the default for the models of \cref{sec:trees}.

\subsection{Local versus Global Explanations}

\subsection{Reading Summary, Force and Dependence Plots}

\subsection{LIME}

%==============================================================================
\section{Data Issues, Bias and Fairness}
\label{sec:data-issues}

\subsection{Imbalanced Datasets and SMOTE}

\subsection{Data Leakage}

\subsection{Fairness Definitions}

%==============================================================================
\section{Experiment Tracking}
\label{sec:experiment-tracking}

\subsection{What to Log and Why}

\subsection{Weights \& Biases and MLflow}
